\section{방법: 증명된 회전과 구성적 확장}
\label{sec:method}

정리~\ref{thm:optimal_hamiltonian}이 직접 정당화하는 것은 \emph{회전 선택}이다: Pre-RoPE 공분산의 PCA 기저로 pre-RoPE 키를 회전시키는 것이 Class~C 내 MSE 기준 최적이다. 그 위의 비트 배분, residual tail 보존, 양자화기 설계는 이 증명 위에 올리는 \emph{구성적 확장}이며, 각각 별도 실험 검증이 필요하다. 본 절은 (i)~증명된 PCA 기반 기본 파이프라인, (ii)~비트 배분 확장, (iii)~per-layer sensitivity 배분을 분리해 기술한다.

\begin{algorithm}[t]
\caption{Pre-RoPE PCA 기반 기본 양자화 파이프라인}
\label{alg:pipeline}
\begin{algorithmic}[1]
\REQUIRE 캘리브레이션 키 $\{k^{(l,h)}_{\mathrm{pre}}\}$ (각 레이어 $l$, KV 헤드 $h$)
\REQUIRE 평균 비트 예산 $B$
\STATE \textbf{[오프라인: 캘리브레이션]}
\FOR{각 레이어 $l$, KV 헤드 $h$}
  \STATE $\Sigma_0^{(l,h)} \leftarrow \frac{1}{n}\sum_i (k_i - \bar{k})(k_i - \bar{k})^\top$ \COMMENT{중심화 공분산}
  \STATE $V_0^{(l,h)}, \{\lambda_j\} \leftarrow \mathrm{eigh}(\Sigma_0^{(l,h)})$ \COMMENT{고유분해}
\ENDFOR
\STATE \textbf{[온라인: 추론 시 k\_proj 후처리]}
\FOR{각 토큰의 pre-RoPE 키 $k_{\mathrm{pre}}$}
  \STATE $z \leftarrow (V_0^{(l,h)})^\top k_{\mathrm{pre}}$ \COMMENT{PCA 좌표 회전}
  \STATE $\hat{z}_j \leftarrow \mathrm{MinMaxQuant}(z_j, B)$ \quad $\forall j$ \COMMENT{기본선: 동일 비트 비대칭 양자화}
  \STATE $\hat{k}_{\mathrm{pre}} \leftarrow V_0^{(l,h)} \hat{z}$ \COMMENT{역회전}
  \STATE RoPE를 $\hat{k}_{\mathrm{pre}}$에 적용 \COMMENT{정상 어텐션 진행}
\ENDFOR
\end{algorithmic}
\end{algorithm}

\paragraph{캘리브레이션.} WikiText-2 train split에서 2048 토큰을 사용하여 각 KV 헤드별로 Pre-RoPE 키 벡터를 추출하고, 중심화 공분산의 고유분해를 수행한다. KVTC와 달리 \emph{헤드별} 독립 PCA를 계산한다 (정리~\ref{thm:optimal_hamiltonian}이 헤드별 최적화를 요구).

\paragraph{양자화.} 기본선은 채널별 비대칭 min-max 균일 양자화다: $\hat{z}_j = \mathrm{round}((z_j - z_{\min}) / s) \cdot s + z_{\min}$, 여기서 $s = (z_{\max} - z_{\min}) / (2^B - 1)$. 이것은 \texttt{k\_proj} 모듈의 forward hook으로 구현되어 모델 코드 수정 없이 적용된다.

\subsection{구성적 확장 1: Water-Filling 비트 배분과 floor 하한}
\label{sec:method_wf}

PCA 회전은 각 차원을 decorrelate하여 독립 스칼라 양자화의 조건을 충족시킨다. 그 위에서 차원별로 다른 비트를 배분하면 추가 개선이 가능하다. 표준 Water-Filling (WF) 배분은 고율 이론에서 유도된다:
\begin{equation}
b_j = B + \frac{1}{2}\log_2\frac{\lambda_j}{\overline{\lambda}}, \qquad b_j \geq b_{\mathrm{floor}},
\label{eq:wf}
\end{equation}
여기서 $\lambda_j$는 PCA 고유값, $\overline{\lambda}$는 기하 평균이다.

\textbf{Floor 하한의 중요성.} 이론적 MSE 최적해는 $b_{\mathrm{floor}} = 0$이다 (Section~\ref{sec:exp_allocation_gap}에서 24/24 케이스로 확인). 그러나 PPL에서는 $b_{\mathrm{floor}} = 2$가 최적이다 (Qwen 2-bit: floor=1 PPL 11.26 $\to$ floor=2 PPL 7.10). 이 ``MSE-최적과 PPL-최적의 괴리''는 Lloyd-Max가 MSE에서 이기지만 PPL에서 실패하는 것과 \emph{동일한 metric mismatch}의 비트 배분 축 발현이다 (Section~\ref{sec:exp_allocation_gap}).

1-bit 양자화는 분산의 63.7\%만 제거하므로 ($c(1) = 0.36$), WF가 1-bit 채널을 생성하면 해당 차원의 양자화 오차가 attention logit을 치명적으로 교란한다. floor$\geq$2를 적용하면 모든 활성 채널이 최소 88.3\%의 분산을 제거하도록 보장한다.

\subsection{Negative result: Query-Weighted Water-Filling (QW-WF)의 무효화}
\label{sec:method_qwwf}

명제~\ref{prop:attn_weighted}에서 attention-weighted distortion은 $\Sigma_Q \cdot \Sigma_K$의 함수이다. 이것은 WF 배분에 query 공분산 정보를 추가하면 개선이 가능할 것이라는 동기를 부여한다:
\begin{equation}
b_j^{\mathrm{QW}} = B + \frac{1}{2}\log_2\frac{\lambda_j^{(K)} \cdot \sigma_j^{(Q)}}{\overline{\lambda^{(K)} \cdot \sigma^{(Q)}}}, \qquad b_j \geq 2,
\label{eq:qwwf}
\end{equation}
여기서 $\sigma_j^{(Q)} = \sqrt{(V_0^\top \Sigma_Q V_0)_{jj}}$는 PCA 좌표에서의 query 가중치다.

\textbf{그러나 이 접근은 실패한다.} 실측에서 QW-WF와 WF(floor=2)의 PPL 차이는 0.5\% 미만이다 (표~\ref{tab:wf_qwwf_honest}). 그 이유는 PCA-Q 자연 정렬 (Section~\ref{sec:exp_pca_q_alignment})에 있다: trained transformer에서 $\lambda_j^{(K)}$와 $\sigma_j^{(Q)2}$의 Spearman 순위 상관이 median $\rho = +0.655$로 강하게 co-monotonic하다. 이 co-monotonicity 하에서 $\sigma_j^{(Q)}$는 $\lambda_j^{(K)}$의 근사적 단조 함수이므로, 식~\eqref{eq:qwwf}의 importance $\lambda_j^{(K)} \cdot \sigma_j^{(Q)}$는 표준 WF의 $\lambda_j^{(K)}$와 동일한 배분 순서를 생성한다. 실측에서 두 배분의 $L^1$ 차이는 budget의 5--8\%이며, 이 5\% bit 변동이 0.5\% PPL 변화로 전이된다.

\textbf{구조적 해석.} Key 고유 기저에서 query 분산의 순위 상관 ($\rho = 0.655$, Section~\ref{sec:exp_pca_q_alignment})으로 인해, query 가중은 기존 WF 배분의 순서를 크게 바꾸지 못한다. 실용적으로는, 비트 배분에 query 정보를 추가하는 것은 불필요하며 표준 WF(floor=2)로 충분하다. 한편, $\Sigma_K$와 $\Sigma_Q$의 고유벡터 부분공간 자체는 30--57°로 정렬되어 있지 않으므로, QW-PCA의 failure는 $\kappa(\Sigma_Q) \sim 10^4$의 수치 불안정에 기인한다.

\subsection{구성적 확장 3: Sensitivity 기반 per-layer 비트 배분}
\label{sec:method_sensitivity}

Section~\ref{sec:exp_per_layer}에서 Lloyd-Max의 PPL 실패가 소수의 레이어에 국소화되어 있음을 확인한다. 이 구조적 발견은 per-layer 비트 배분이라는 직접적 처방을 제공한다:

\begin{enumerate}[leftmargin=1.5em,itemsep=2pt]
\item 각 레이어를 개별적으로 $b$-bit으로 양자화하고, 나머지를 FP16으로 유지하여 $\Delta\mathrm{PPL}_l$를 측정한다.
\item $\Delta\mathrm{PPL}_l$가 큰 상위 $k$개 레이어에 $B_{\mathrm{high}}$-bit을, 나머지에 $B_{\mathrm{low}}$-bit을 배분한다.
\end{enumerate}

이것은 차원별 WF (식~\ref{eq:wf})와 직교하는 \emph{레이어별} 배분이다. 두 배분을 결합하면 (sensitive layer에 3-bit + 차원별 WF) 추가 개선이 가능하지만, 본 논문에서는 각 축을 독립적으로 평가한다.

\paragraph{계산 비용.} 캘리브레이션은 $O(n d_{\mathrm{head}}^2)$ (고유분해 1회), 추론 시 $O(d_{\mathrm{head}}^2)$ (행렬-벡터 곱 2회)이다. Sensitivity 측정은 $O(L)$회의 단일-레이어 양자화 PPL 평가가 추가로 필요하다 (Mistral 32 layers: $\sim$70초).
